% Part: sets-functions-relations
% Chapter: infinite
% Section: dedekind
%
\documentclass[../../../include/open-logic-section]{subfiles}

\begin{document}

\olfileid[id]{sfr}{infinite}{dedekind}
\olsection{Aljabar Dedekind}

Kita tidak hanya menghendaki bilangan asli itu tak hingga; kita menghendaki
agar bilangan-bilangan tersebut mempunyai sifat-sifat (aljabar) tertentu:
bilangan-bilangan itu harus berperilaku baik terhadap penjumlahan, perkalian,
dan seterusnya.

Gagasan Dedekind ialah menjadikan gagasan \emph{fungsi penerus} sebagai
dasar, lalu mencirikan bilangan sebagai objek yang mempunyai sifat-sifat
berikut:
\begin{enumerate}
	\item Terdapat suatu bilangan, $0$, yang bukan penerus bilangan mana pun
	\\yakni, $0 \notin \ran{s}$
	\\yakni, $\forall x\ s(x) \neq 0$
	\item Bilangan-bilangan yang berbeda mempunyai penerus yang berbeda
	\\yakni, $s$ adalah !!a{injection}
	\\yakni, $\forall x \forall y (s(x) = s(y) \lif x = y)$
	\item\ollabel{repeatedapplication} Setiap bilangan diperoleh dari
	$0$ melalui penerapan berulang fungsi penerus.
\end{enumerate}
Dua syarat pertama mudah ditangani dengan logika orde pertama (lihat di
atas). Namun, kita tidak dapat menangani \olref{repeatedapplication} hanya
dengan logika orde pertama. Terobosan Dedekind ialah merumuskan ulang syarat
\olref{repeatedapplication}, dalam istilah teori himpunan, sebagai berikut:
\begin{enumerate}
	\item[3$'$.] Bilangan asli adalah himpunan terkecil yang
	\emph{tertutup terhadap fungsi penerus}: yakni, jika kita menerapkan
	$s$ pada sembarang !!{element} dari himpunan tersebut, kita memperoleh
	!!{element} lain dari himpunan itu.
\end{enumerate}
Namun, kita perlu menguraikan hal ini secara perlahan.

\begin{defn}\ollabel{Closure}
	Untuk sembarang fungsi $f$, himpunan $X$ bersifat $f$-\emph{tertutup}
	{jika dan hanya jika} $(\forall x \in X)f(x) \in X$. Sekarang, untuk setiap $o$,
	definisikan
	$$\closureofunder{f}{o} = \bigcap\Setabs{X}{o \in X\text{ dan }X
\text{ bersifat $f$-tertutup}}$$
\end{defn}

Jadi, $\closureofunder{f}{o}$ adalah irisan semua himpunan $f$-tertutup
yang mempunyai $o$ sebagai !!a{element}. Secara intuitif,
$\closureofunder{f}{o}$ merupakan himpunan $f$-tertutup \emph{terkecil}
yang mempunyai $o$ sebagai !!a{element}. Hasil berikutnya membuat gagasan
intuitif itu menjadi saksama:
\begin{lem}\ollabel{closureproperties}
	Untuk sembarang fungsi $f$ dan sembarang $o \in A$:
	\begin{enumerate}
		\item\ollabel{closurehaselem} $o \in \closureofunder{f}{o}$; dan
		\item\ollabel{closureclosed} $\closureofunder{f}{o}$ bersifat $f$-tertutup; dan
		\item\ollabel{closuresmallest} jika $X$ bersifat $f$-tertutup dan $o \in
		X$, maka $\closureofunder{f}{o} \subseteq X$.
	\end{enumerate}
\end{lem}

\begin{proof}
Ada sekurang-kurangnya satu himpunan $f$-tertutup
yang mempunyai $o$ sebagai !!a{element}, yakni $\ran{f}\cup \{o\}$. Jadi,
$\closureofunder{f}{o}$, yaitu irisan dari \emph{semua} himpunan semacam itu,
ada. Sekarang kita harus memeriksa
\olref{closurehaselem}--\olref{closuresmallest}.

Mengenai \olref{closurehaselem}: $o \in \closureofunder{f}{o}$ karena
himpunan tersebut merupakan irisan himpunan-himpunan yang semuanya mempunyai
$o$ sebagai !!a{element}.

Mengenai \olref{closureclosed}: andaikan $x \in \closureofunder{f}{o}$.
Jadi, jika $o \in X$ dan $X$ bersifat $f$-tertutup, maka $x \in X$; dan
karena $X$ bersifat $f$-tertutup, sekarang $f(x) \in X$. Maka
$f(x) \in \closureofunder{f}{o}$.

Mengenai \olref{closuresmallest}: secara umum, jika $X
\in C$, maka $\bigcap C \subseteq X$.
\end{proof}

Dengan ini, kita dapat mengatakan:

\begin{defn}
Suatu \emph{aljabar Dedekind} adalah suatu himpunan $A$, bersama dengan
fungsi $f \colon A \to A$ dan suatu $o \in A$, sedemikian sehingga:
	\begin{enumerate}
		\item \ollabel{ded:proper} $o \notin \ran{f}$
		\item \ollabel{ded:injection} $f$ adalah !!a{injection}
		\item \ollabel{ded:closure} $A = \closureofunder{f}{o}$
	\end{enumerate}
\end{defn}

Karena $A = \closureofunder{f}{o}$, hasil kita sebelumnya memberi tahu bahwa
$A$ adalah himpunan $f$-tertutup terkecil yang mempunyai $o$ sebagai
!!a{element}. Jelaslah bahwa suatu aljabar Dedekind bersifat tak hingga
Dedekind; lihat saja klausa \olref{ded:proper} dan
\olref{ded:injection} dalam definisi tersebut. Namun, fakta yang lebih
menarik ialah bahwa setiap himpunan yang tak hingga Dedekind dapat dijadikan
suatu aljabar Dedekind.

\begin{thm}\ollabel{thm:DedekindInfiniteAlgebra}
Jika terdapat suatu himpunan yang tak hingga Dedekind, maka terdapat suatu
aljabar Dedekind.
\end{thm}

\begin{proof}
Misalkan $D$ bersifat tak hingga Dedekind. Jadi, terdapat suatu injeksi
$g \colon D \to D$ dan suatu elemen $o \in D \setminus \ran{g}$.
Sekarang, misalkan $A = \closureofunder{g}{o}$; berdasarkan
\olref{closureproperties}, $A$ ada dan $o \in A$. Misalkan
$f = \funrestrictionto{g}{A}$. Kita akan menunjukkan bahwa $A, f, o$
membentuk suatu aljabar Dedekind.

Mengenai \olref{ded:proper}: $o \notin \ran{g}$ dan $\ran{f}
\subseteq \ran{g}$, sehingga $o\notin \ran{f}$.

Mengenai \olref{ded:injection}: $g$ adalah suatu injeksi pada $D$; maka
$f \subseteq g$ harus merupakan suatu injeksi.

Mengenai \olref{ded:closure}: berdasarkan \olref{closureproperties},
$A$ bersifat $g$-tertutup; terlebih lagi, $A$ bersifat $f$-tertutup. Maka
$\closureofunder{f}{o} \subseteq A$ berdasarkan \olref{closureproperties}.
Karena $\closureofunder{f}{o}$ juga bersifat $f$-tertutup dan
$f = \funrestrictionto{g}{A}$, diperoleh bahwa $\closureofunder{f}{o}$
bersifat $g$-tertutup. Jadi, $A \subseteq \closureofunder{f}{o}$ berdasarkan
\olref{closureproperties}.
\end{proof}

\end{document}
